\section{研究方向概述}
\label{section-generic}

控制系统是现代工业和社会基础设施的核心组成部分，其功能安全与可靠性对自动化生产、智能交通以及能源管理等关键领域的运行效能与整体安全至关重要。这类系统对实时响应、复杂交互及高可靠运行要求非常严格。功能安全的保障贯穿其整个生命周期，需要通过严谨的设计、开发和验证流程，确保系统在正常甚至故障情况下，仍能有效执行关键的安全功能，从而将风险降到最低。在控制系统功能安全保障中，需求分析阶段的高质量建模具有基础性作用，它为后续的设计实现和验证提供了必要的依据。鉴于控制系统本质上是状态驱动的动态系统，状态机建模因其能清晰地表达系统离散运行模式、状态转移逻辑及事件驱动的行为，成为描述其复杂动态行为和刻画安全约束的主流形式化方法。统一建模语言（UML）状态图支持层次化状态、并发、守卫条件、动作等丰富语义，因此被广泛用于精确描述具有高度交互性和实时性要求的复杂控制系统行为。

然而，随着控制系统规模和复杂性的持续增加，传统的状态机建模方法面临着显著的瓶颈。首先，建模过程非常依赖专家经验并且成本很高。工程师需要从表述抽象、来源多样的自然语言需求文档中，精确地提取关键要素，然后将其准确地映射为符合形式化语义的模型元素。这是一个耗时、易出错且非常耗费人力的工作。早期基于规则的自动化建模方法旨在减轻人工负担，但其规则库的构建和维护本身就非常复杂，难以灵活适应需求的动态变化并充分融入多样化的领域知识。其次，模型可信度的验证也面临挑战。状态机模型的正确性与完备性对系统安全分析至关重要。形式化验证方法虽然理论上很完备，但在处理复杂系统时容易遇到状态空间爆炸的问题，并且对工程师的形式化方法素养要求很高，在工程实用性和普适性方面受到限制。基于测试的验证方法在工程上更友好，但其效果高度依赖于测试场景的质量和覆盖程度。人工设计能充分覆盖关键路径、边界条件及异常情况的测试场景同样需要大量精力并依赖专家经验，效率低下且难以保证测试的充分性。

近年来，大语言模型（Large Language Models, LLM）在深度理解自然语言和生成结构化内容方面展现出卓越的能力，这为自动化的需求分析与建模开辟了新的途径。研究表明，LLM能有效地解析需求文档语义，识别关键实体及其关系 \upcite{Vega2024LLMCapabilities}，并具备按特定规则生成初步模型描述或架构的潜力 \upcite{Wei2024Requirements}。在验证环节，利用LLM辅助生成多样化的测试场景也显示出初步的可行性 \upcite{Khan2025LLMRequirements}，这为解决建模和验证场景生成的高成本问题提供了潜在的技术机遇。

然而，将LLM应用于控制系统状态机建模与验证仍然存在关键挑战。在建模方面，现有的基于LLM的单步方法在处理复杂需求时，常常缺乏对状态机固有结构特征的有效引导和对领域知识的系统性融合，导致生成的模型在完整性、语义准确性及领域约束符合性方面存在不足。在验证方面，虽然LLM在生成描述性测试场景上展现出潜力，但针对功能安全的核心目标，仅在行为层面进行仿真测试可能不够充分且效率有限。保障功能安全通常要求模型满足特定的、可形式化表述的关键性质，因此迫切需要探索如何利用LLM高效地生成待验证的系统性质以及用于检验此类核心性质的验证激励（例如用于形式化验证的输入序列）。同时，还需要设计能够有效执行此类性质验证、尤其能克服传统形式化方法状态爆炸瓶颈的高效验证机制。此外，在模型修复环节，现有的基于LLM的自我反思机制大多未能有效利用验证环节产生的具体反馈信息，例如性质违背后的反例，限制了基于反馈进行迭代修复的效率和可靠性。

综上所述，高效、准确地针对控制系统功能需求进行状态机建模，是保障系统功能安全的基石，但传统方法在应对复杂性和提高效率方面存在根本的局限。LLM为解决自动化建模与验证提供了新的机遇，但在状态机建模的引导、面向形式化性质的高效验证方法创新，以及基于验证反馈（特别是利用反例）的迭代优化等方面仍存在显著挑战。因此，深入研究如何有效融合LLM的强大语义处理能力、状态机的结构化特征、领域知识以及带有反例指导的验证反馈机制，建立一套迭代式的模型构建、性质验证与修复方法，具有重要的理论价值和广泛的工程应用前景。本文将从状态机建模方法、自动化模型生成、状态机缺陷检测与形式化验证以及模型修复方法四个方面，系统梳理国内外研究现状，为后续研究提供理论基础与技术参照。